\chapter{From SF Research to the Integrated Barrick Project}
\chapterstatus{The project history in this chapter follows the mandate supplied
by SF Research. The methodological map is a source-grounded synthesis of Teams
1--8 and the accepted current evidence; it does not promote the conditional
corporate proxy to a filing-reconciled valuation.}

\section{Research question, community and valuation object}

% Sources: CONTENT-LEDGER BG-U-01, BG-U-10--14; CL-001/002/011/013;
% Team 8 T8-17--19; Team 4 MAP-T4-08--12; Team 5 T5-02--19;
% SF Research complete guide, project organisation and unification workflow.
SF Research did not begin as a preassembled quantitative-finance laboratory.
It originated as a comparatively vertical group, strongly represented by
management-engineering students and initially centred on events and financial
education. Research activity first appeared through short studies developed by
small groups over roughly monthly horizons. As contributors with different
academic paths entered the area, those isolated exercises gradually revealed a
larger opportunity: to build a common annual project in which every member
could contribute while learning how separate disciplines meet in a real
valuation problem.

That evolution matters because the project's main advantage is
interdisciplinarity. A mathematics curriculum can reasonably treat probability,
analysis or numerical foundations as prerequisites; a management curriculum
can do the same with accounting and valuation; a computing curriculum can
assume familiarity with implementation. The team deliberately avoided those
silent assumptions. The project was designed from the bottom up so that it
could extend, rather than merely repeat, the contributors' university
programmes and introduce models and techniques that many of them had not met
in formal courses.

The resulting chain joins economics and corporate finance with mathematics,
econometrics, statistics and probability, stochastic methods, numerical
methods, optimisation, operations research and computing. At university these
subjects are often taught as distinct modules. In the Barrick project they are
forced to interact: probability supplies the language of uncertainty;
econometrics tests the historical series; stochastic calculus defines the
continuous-time models; numerical methods make calibration and simulation
possible; optimisation estimates parameters; operational modelling converts
gold, production and costs into margins; and corporate finance closes the DCF.
The thesis is therefore both a valuation study and a record of how theoretical
ideas become mutually dependent when applied to one object.

The integrated research question is: \emph{how can uncertainty in commodity
prices, production, costs and corporate assumptions be propagated into a
distribution of Barrick free cash flow and value without confusing
risk-neutral option information with a physical operating forecast?} The
object is not a point price produced by an isolated option model. It is a
conditional distribution of enterprise value and, after a separately verified
capital bridge, equity value and value per diluted share. The comparison with
the observed market price is therefore a dated diagnostic inside a declared
assumption set, not a target price or a recommendation.

This project was developed within Research, the Starting Finance Club PoliTo
area dedicated to collaborative quantitative, economic, financial and
methodological projects. Its members bring heterogeneous backgrounds---from
management engineering and mathematics to data science, computational methods
and quantitative finance---so the research process is organised around
knowledge sharing rather than an assumed common starting point. The practical
objective is to establish enough shared language and method for contributors
with different training to make a verifiable contribution, while keeping
technical depth accessible to readers who did not participate in the project.
This description concerns how the work was organised; it is not presented as a
formal institutional mission or organisational chart.

That collaborative setting is material to the valuation design. The commodity
block must state whether paths are risk-neutral, physical or explicit
scenarios; the operating block must state production, cost and realised-price
units; the accounting block must reconcile operating margin, depreciation,
taxes and reinvestment; and the valuation block must state discounting,
terminal closure and every enterprise-to-equity adjustment. A result becomes
economically interpretable only when contributors working on different modules
use the same date, perimeter, currency, notation and manifest. The research
question is therefore also an integration question: whether heterogeneous
analytical work can be connected without hiding disagreements in definitions
or silently filling missing data.

\section{From eight autonomous studies to one chain of evidence}

% Sources: CONTENT-LEDGER authority table; MAP-T1--T4; T5-01--T8-20;
% SF Research complete guide, unification and review workflow.
The eight source studies divide the problem into complementary modules. Team 1
provides probability, measure and Black--Scholes foundations; Team 2 provides
regression and linear time-series methods; Team 3 provides Monte Carlo and
simulation diagnostics; Team 4 maps production, cost and price scenarios into
an operating layer; Team 5 supplies DCF, FCFF, WACC and stochastic-valuation
concepts; Team 6 develops stochastic variance, jumps and Fourier pricing; Team
7 supplies historical gold, dependence and volatility evidence; and Team 8
compares option-surface models and produces risk-neutral GLD paths. These were
autonomous studies with different immediate questions, conventions and
validation states, not chapters originally written as a single sequential
paper.

Unification was consequently treated as a separate research phase rather than
as copy-and-paste assembly. Before drafting the integrated argument, the work
froze the admissible sources and material claims, assigned an authority to each
concept, and classified outputs as theoretical, synthetic, legacy empirical,
current empirical or newly generated. Data and notation contracts then made
dates, currencies, units, measures and interfaces explicit. Source-model
implementations were preserved and subjected to parity tests before being
connected to the unified runner, while the current LSE refresh was recorded as
a separate evidence layer so that a new market observation could not be
mistaken for a re-estimation of every legacy study.

The execution itself followed two connected workstreams. The code stream built
and tested the valuation pipeline, generated tables and figures from accepted
manifests, and exposed unresolved inputs instead of replacing them with silent
assumptions. The writing and LaTeX stream harmonised notation, removed
duplication, created transitions and kept claims tied to their evidence status.
Technical review checked models, calculations and reproducibility; editorial
review checked scope, language and reader navigation; final quality assurance
tested both the compiled artifact and the consistency between prose, tables,
figures and generated outputs. Section~\ref{sec:team-study-synthesis} makes the
resulting roles and validation states directly comparable.

\section{Contribution, falsifiable boundaries and roadmap}

% Sources: VAL-001 gate sequence in GAP-REGISTER; CL-001/005/012/013.
The first contribution is an auditable handoff from option-implied dynamics to
corporate quantities. The second is an evidence hierarchy that separates
source-model parity, market-data validity, statistical fit, predictive tests,
numerical convergence and accounting reconciliation. The third is a reporting
contract: a stochastic valuation must expose its median and tails, the
contemporaneous market-price reference and the assumptions that move the
distribution, rather than presenting a single estimate with false precision.
Together, these contributions make disagreement inspectable: a reader can see
whether a result changes because of market evidence, model form, operating
assumptions, accounting treatment or the equity bridge.

The same structure creates falsifiable boundaries. Better option fit does not
prove better out-of-sample forecasts; better forecasts do not prove a valid
GLD-to-realised-gold bridge; a valid commodity bridge does not prove the mine
or accounting perimeter; and a coherent enterprise value does not determine
equity value until net debt, minorities, non-operating assets and diluted
shares are reconciled. Each boundary has a corresponding validation gate in
Chapter 12. Failure at a later gate does not erase an earlier technical
result, but it narrows the conclusion that the integrated study may draw from
that result.

The thesis follows the same chain. Part I reconstructs the genesis of the
research, the monetary and macroeconomic role of gold, and the Barrick
valuation perimeter. Part II consolidates the mathematical, econometric and
numerical prerequisites from the bottom up. Part III first tests the legacy
gold series with discrete-time methods, then introduces the affine
stochastic-volatility and jump model ladder, and finally calibrates and
compares the models on the option surface. Part IV carries the admitted gold,
production and cost paths through EBITDA and DCF, validates the integrated
experiment and closes with limitations and future research. No part of the
study provides a buy, sell or hold recommendation.
